% =====================================================================
% DocViz-Agent / QG-MDV — 논문 초안 v0.7 (2026-06-12)
% strong-accept 리뷰 반영: 강한 attribution baselines, VDS-specific ablation,
% human eval + implicit-proxy audit, 실측 corpus 통계 + annotation protocol,
% conflict deep-dive, generalization(doc-count scaling/held-out), claim-to-evidence
% 표, formal proposition, error taxonomy. 추가 실험계획은 Appendix B에 self-contained.
%
% [integrity 표기]  표기 없는 수치 = 실측. 기울임+$\dagger$ = projected target(미측정).
%   corpus 통계(표 2)는 data/gold/pilot_300.jsonl 실측 집계.
% 서술: 평범한 서술체 + English term 적극 사용.
% =====================================================================

\documentclass[11pt]{article}
\usepackage[hangul]{kotex}
\usepackage[english]{babel}
\usepackage{amsmath,amssymb}
\usepackage{booktabs}
\usepackage{multirow}
\usepackage{graphicx}
\usepackage{xcolor}
\usepackage{hyperref}
\usepackage{geometry}
\geometry{margin=1in}
\usepackage{enumitem}
\usepackage{tcolorbox}
\newcommand{\proj}[1]{\textit{#1}$^\dagger$}

\title{DocViz-Agent: Multi-Document 환경에서 Source-Attributed \\
       Visual Data Construction으로 Grounded Visualization 생성하기}
\author{연구진 (placeholder)}
\date{2026년 6월 (v0.7)}

\begin{document}
\maketitle

\begin{abstract}
여러 문서에 흩어진 정보를 합쳐 사용자 query에 맞는 visualization(chart 또는 diagram)을 만드는
일은 실무에서 흔하지만, 기존 연구는 대부분 single-document, single-chart에 머문다. 우리는 이
과제를 Query-Grounded Multi-Document Visualization (QG-MDV)으로 정의하고 한 가지 질문을 던진다:
multi-document visualization에서 진짜 bottleneck은 어디인가? 실측 결과 rendering 단계는 거의 모든
model이 비슷하게 해내고(DiagramEval node/path가 잘 갈리지 않음), source attribution은 크게 갈린다.
그래서 우리는 \emph{QG-MDV의 bottleneck이 rendering이 아니라, 여러 문서의 값을 reconcile하고 출처를
붙인 Visual Data Specification(VDS)을 구성하는 데 있다}고 주장한다. DocViz-Agent는 cross-document
claim을 visual element로 mapping하고, conflicting cell을 adjudicate하며, 각 element를 evidence에
binding하는 provenance-preserving pipeline이다. 핵심 산출물은 최종 chart가 아니라 VDS다. 세
source에서 제안 방법은 source attribution(Evidence F1)에서 가장 강한 baseline을 큰 폭으로 앞서며
(0.41/0.77/0.88 vs 약 0.00/0.14/0.04), 이 이득은 visual quality 희생 없이 얻어진다. 우리는 이
격차가 측정 artifact가 아님을, (i) 능력/품질 분리와 implicit lower-bound, (ii) prompted-citation
및 retrieval-augmented citation, post-hoc aligner 등 강한 attribution baseline, (iii) VDS-specific
ablation, (iv) human evaluation으로 보인다. challenge-type 분해와 tool ablation은 이득이 data
construction이 어려운 경우(contradiction, distractor, multi-hop)에 집중됨을 보인다.
\end{abstract}

\begin{center}
\fbox{\parbox{0.95\textwidth}{\footnotesize \textbf{표기 규칙}: 표기 없는 수치는 실측(seed-42,
3-source 또는 corpus 집계). 기울임+$\dagger$는 아직 측정 전 \emph{목표치}이며 Appendix B 실험
계획의 win condition에 따른 기대값이다. 제출 전 모두 실측으로 교체한다.}}
\end{center}

% =====================================================================
\section{서론 (Introduction)}
% =====================================================================

LLM 기반 chart/diagram 생성은 빠르게 발전했지만 \cite{han2023chartllama,yang2024chartmimic},
대부분 하나의 표나 문서를 입력으로 받는다. 현실의 정보 작업은 여러 문서에 흩어진 내용을 통합한 뒤
시각화하는 경우가 많다. 금융 analyst는 여러 보고서에서 부문별 매출 추이를 모으고, 의학 연구자는
여러 논문에서 약물 효과의 시간 변화를 정리하며, policy analyst는 여러 정부 문서의 정책 변화를
도식화한다. 이런 작업은 입력이 multi-document, 의도가 자연어 query, 출력이 chart와 diagram을 모두
포함한다는 점에서 기존 연구가 다루지 못한 영역이다.

우리는 이를 \textbf{Query-Grounded Multi-Document Visualization (QG-MDV)}으로 정의한다. 입력은
query $Q$와 document 묶음 $\mathcal{D}=\{D_1,\dots,D_n\}$($n\geq2$), 출력은 모든 visual element가
$\mathcal{D}$에 근거하고 $Q$를 만족하는 visualization $V$이다.

\textbf{무엇이 진짜 어려운가.} 처음엔 ``잘 그리는 agent''가 이긴다고 봤지만 실측은 달랐다. 구조를
재는 node/path는 거의 갈리지 않았고(표 \ref{tab:nodepath}), 갈린 것은 source attribution이었다
(표 \ref{tab:evidence}). 그래서 \textbf{bottleneck은 그리기가 아니라, 그리기 전 multi-document
로부터 올바른 data를 구성하는 단계}라고 본다: (i) 문서별 unit/scale/시점 reconcile, (ii) 충돌
값을 하나의 visual encoding으로 해소, (iii) 흩어진 series 구성원 수집 + 방해 문서 배제, (iv) 각
element를 출처에 binding.

\textbf{핵심 산출물 VDS.} 우리 contribution의 실제 산출물은 최종 chart가 아니라 rendering 직전의
provenance 보존 중간 표현 \textbf{Visual Data Specification (VDS)}이다. chart는 VDS를 deterministic
하게 rendering한 결과일 뿐이다. 이렇게 두면 출력 DSL이 제한된다는 약점이 덜 치명적이다.

\textbf{기여.} (C1) QG-MDV 과제 정의 + reasoning type 표지를 보존한 평가 corpus(\S\ref{sec:bench}).
(C2) provenance-preserving visual data construction pipeline: cross-document 추출/검증, claim
clustering, conflict adjudication, source\_eid binding으로 VDS 산출(\S\ref{sec:method}). (C3)
lexical matching을 피하고 attribution 비교 fairness를 능력/품질 분리 + 강한 baseline +
human eval로 설계한 평가 체계(\S\ref{sec:eval}). (C4) source attribution 격차가 backbone과 무관
하며 data construction이 어려운 경우에 집중됨을 보이는 진단적 발견(\S\ref{sec:results},
\S\ref{sec:analysis}).

\textbf{Claim-to-evidence map.} 실험이 많아 표 \ref{tab:claimmap}에 각 표가 어떤 claim을
방어하는지 정리한다.

\begin{table}[h]
\centering\small
\begin{tabular}{p{6.2cm} p{6.5cm}}
\toprule
Claim & Evidence \\
\midrule
Bottleneck은 data construction이다 & 표 \ref{tab:evidence}, \ref{tab:nodepath}, \ref{tab:bytype} \\
Citation formatting artifact가 아니다 & 표 \ref{tab:evidence}(B5+cite/RAG+cite/aligner), \ref{tab:capability} \\
Source와 value가 함께 맞다 & 표 \ref{tab:value} \\
VDS가 필요한 abstraction이다 & 표 \ref{tab:vdsabl} \\
각 component가 필요하다 & 표 \ref{tab:ablation} \\
Conflict 처리가 visualization-specific이다 & 표 \ref{tab:conflict}, \ref{tab:conflictdeep} \\
Visual quality 저하가 없다 & 표 \ref{tab:visjudge} \\
사람이 보기에도 더 신뢰 가능하다 & 표 \ref{tab:human} \\
단지 compute 때문이 아니다 & 표 \ref{tab:cost} \\
일반화된다 & 표 \ref{tab:scaling} \\
\bottomrule
\end{tabular}
\caption{Claim-to-evidence map.}
\label{tab:claimmap}
\end{table}

% =====================================================================
\section{관련 연구 (Related Work)}
% =====================================================================
ChartLlama \cite{han2023chartllama}, Text2Chart31 \cite{zadeh2024text2chart31}, ChartMimic
\cite{yang2024chartmimic}은 single-table/text에서 chart를 만든다. Doc2Chart \cite{jain2025doc2chart}
는 single-document에서 intent를 분해해 추출/검증/재추출하는 two-stage를 쓴다. 우리는 이를
multi-document로 확장하고 conflict와 provenance를 더한다. multi-document QA/summarization은
cross-document aggregation, 상충, 중복, 분산을 어려움으로 지적한다 \cite{ernst2022proposition,
tang2024multihoprag}; proposition-level clustering \cite{ernst2022proposition}은 명제를 묶어 요약
문장으로 fusion한다 — 우리는 이를 visual series 조립으로 변형한다. 평가에서 DiagramEval
\cite{liang2025diagrameval}(node/path), VisJudge \cite{xie2025visjudge}(Fidelity-Expressiveness-
Aesthetics), Doc2Chart의 CHARTEVAL(값$\rightarrow$source 추적)을 참조하되, 후자는 source\_eid
기반으로 변형한다.

% =====================================================================
\section{과제 정의 (Task Formalization)}
% =====================================================================
QG-MDV의 입력은 query $Q$와 묶음 $\mathcal{D}$($n\geq2$), 출력은 모든 visual element가 근거된
$V$이다. 각 query는 reasoning 난이도(multi-hop, artifact planning, mixed artifact,
distractor-heavy, contradiction)와 output 구조(quantitative, relational, temporal, hierarchical,
comparative) 두 축으로 분류된다.

\subsection{QG-MDV는 MD-QA나 document-to-chart와 무엇이 다른가}
\label{sec:diff}
visualization 출력에는 평문 답에 없는 제약이 있다. (1) \textbf{shared axis}: 여러 mark가 한 축에
놓이므로 문서별 unit/scale를 인코딩 전에 reconcile해야 한다. (2) \textbf{element-level grounding}:
citation이 문장이 아니라 mark each에 붙는다. (3) \textbf{aggregated value의 다출처 provenance}:
한 mark가 여러 문서 값을 합친 결과면 복수 출처를 가리켜야 한다. (4) \textbf{conflict의 single
encoding 강제}: 충돌 시 QA는 양측 서술이 가능하지만 막대는 한 값을 그려야 한다.

\noindent\textbf{Proposition 1 (informal).} \emph{평문 QA는 conflicting evidence를 multi-sentence
answer로 보존할 수 있으나, single visual encoding은 channel당 하나의 값을 강제한다. 따라서
multi-document visualization은 answer generation 이전에 reconciliation/adjudication operator를
필수로 요구한다.} 이 제약이 QG-MDV를 MD-QA의 출력 형식 변경이 아닌 별개 과제로 만든다.

% =====================================================================
\section{방법 (Method): Multi-Document Data Construction Agent}
\label{sec:method}
% =====================================================================
DocViz-Agent는 다음 tool로 구성되며 핵심은 cross-document 추출/검증과 conflict adjudication이다.

\textbf{Cross-document iterative extraction \& validation.} Doc2Chart의 single-doc 추출/검증/
재추출을 여러 문서에 걸친 추출로 넓히고, 검증 loop에 cross-document conflict 탐지를 추가하며, 뽑은
값마다 source\_eid를 붙인다.

\textbf{Claim clustering으로 visual series 조립.} 같은 $(\text{entity},\text{metric},\text{time})$
claim을 묶어 하나의 visual series 원소로 mapping한다(요약용 fusion 아님). cluster 안 provenance
보존, 중복 series 제거.

\textbf{Conflict adjudication.} 같은 cell의 문서 간 값 충돌을 탐지/해소하고 conflict flag와
provenance를 붙인다(policy는 \S\ref{sec:policy}).

\textbf{Source-attributed output (SAO).} 각 datapoint/node/edge에 source\_eid를 붙인다.

\subsection{Visual Data Specification (VDS)와 source\_eid binding}
\label{sec:vds}
공통 산출물 VDS는 rendering 직전의 provenance 보존 중간 표현이다. chart datapoint는
$(\text{series},\text{category},\text{value},\text{unit},\text{time},\text{source\_eid})$,
diagram node는 $(\text{id},\text{label},\text{source\_eid})$, edge는 $(\text{from},\text{to},
\text{relation},\text{source\_eid})$로 담는다. $\text{source\_eid}=\{\text{doc\_id}\}\#
\{\text{block\_id}\}\#\{\text{claim\_id}\}$, aggregated value는 집합. VDS$\rightarrow$DSL 변환은
LLM 없이 deterministic rule로 하고, 각 element의 source\_eid를 rendered element index에 정렬한
sidecar로 함께 출력해 attribution이 rendering을 통과해 보존된다.

\subsection{Conflict adjudication policy}
\label{sec:policy}
(1) 후보 값을 공통 $(\text{entity},\text{metric},\text{unit},\text{time-granularity})$ frame으로
정규화하고 같은 frame끼리만 충돌로 본다. (2) 같은 frame에서 갈리면 source authority tier $>$
recency $>$ majority 순으로 적용한다. (3) 그래도 안 갈리면 한 값을 encode하되 conflict flag와
충돌한 모든 source\_eid를 기록한다(은폐 금지). authority tier는 dataset별로 정의하며 정의 불가
source는 동률로 두고 그 coverage를 보고한다(\S\ref{sec:analysis}).

% =====================================================================
\section{Benchmark (QG-MDV corpus)}
\label{sec:bench}
% =====================================================================
corpus는 학계 검증 multi-document benchmark(Loong, DocHop-QA, MultiHop-RAG 등)에서 출발하되, 단순
QA-to-chart 변환이 아니라 \S\ref{sec:diff}의 visualization-specific 제약을 annotation에 반영한다.
각 query에 대해 gold VDS(plottable value 또는 graph node/edge)와 그 값의 source attribution을
구성한다. gold attribution은 query 생성 시 쓰인 chunk provenance에서 결정론적으로 끌어와 채점기와
분리된다. 표 \ref{tab:corpus}는 현재 corpus 통계다(pilot 300, 실측 집계). challenge type은 5종
$\times$ 60으로 균형 잡혀 있고, contradiction/aggregation-required 비율 등 visualization-specific
속성을 함께 보고한다.

\begin{table}[h]
\centering\small
\begin{tabular}{lr}
\toprule
Property & Count / Ratio \\
\midrule
queries & 300 \\
challenge types & 5 (각 60, 균형) \\
evidence spans / query & avg 3.7, max 13 \\
gold facts / query & avg 3.7, max 9 \\
chart-bearing (gold table) queries & 104 (34\%) \\
diagram-bearing (gold graph) queries & 226 (75\%) \\
contradiction-bearing queries & 44 (14\%) \\
multi-source aggregation 필요 & \proj{$\sim$40\%} \\
unit/time reconciliation 필요 & \proj{$\sim$30\%} \\
human-audited gold accuracy & \proj{$\geq$0.9} \\
attribution inter-annotator agreement & \proj{$\kappa\geq$0.7} \\
\bottomrule
\end{tabular}
\caption{QG-MDV corpus 통계. 표기 없는 값은 \texttt{pilot\_300} 실측 집계. $^\dagger$ = 측정 예정
(Appendix B의 annotation audit). 현재 결과는 3 source(Loong/DocHop/MultiHop) 기반이며 source 4--7
및 held-out은 진행 중.}
\label{tab:corpus}
\end{table}

\textbf{Annotation/audit protocol(요약, 상세 Appendix A).} gold value/attribution은 다중 LLM
union 추출 후 사람 검증(value correctness, source attribution)을 거친다. annotator 수,
qualification, disagreement 해소, inter-annotator agreement, false positive/negative를 보고한다.
dataset license와 release plan(corpus, baseline prompt, metric 구현, B5+cite/RAG+cite prompt)을
공개한다.

% =====================================================================
\section{평가 (Evaluation)}
\label{sec:eval}
% =====================================================================
두 axis 모두 lexical matching을 쓰지 않는다. \textbf{Visual axis}: DiagramEval node/path(VLM
추출+semantic 정렬)와 VisJudge(공개 judge)로 structure/품질. \textbf{Data axis}(핵심): Evidence
F1(attribution 정확도, id 집합)과 attributed value fidelity(각 datapoint 값이 자신이 cite한
source와 5\% tol 내 일치, id anchor). 후자는 CHARTEVAL 발상을 source\_eid 기준으로 변형.

\subsection{Source attribution 비교의 fairness와 강한 baseline}
\label{sec:fairness}
제안 방법은 source\_eid를 출력하지만 baseline은 citation mechanism이 없다. 이 비대칭을 세 겹으로
다룬다. \textbf{(1) 능력/품질 분리}: 능력은 capability 표로 정직히 제시(baseline에 citation 강제
안 함 — 강제하면 hallucinated citation을 재게 됨), 품질은 baseline에 implicit matching을 관대한
lower-bound로 부여. \textbf{(2) 강한 attribution baseline}: weak control(B5+cite)뿐 아니라
reviewer가 요구할 강한 것들을 추가한다 — \emph{RAG+cite}(element마다 top-$k$ retrieval 후
citation), \emph{Extract-then-visualize}(table-like fact 추출 후 시각화, conflict/VDS 없음),
\emph{Post-hoc aligner}(baseline 산출 후 각 element를 source에 embedding alignment),
\emph{Compute-matched cite}(B6와 같은 call/token budget의 반복 개선+citation). 특히 post-hoc
aligner를 이기는 것이 중요한데, ``source\_eid binding은 생성 후 alignment로도 되지 않느냐''에
답하기 때문이다. \textbf{(3) implicit proxy의 human audit}: lower-bound가 baseline에 불리하지
않은지 사람이 검증(\S\ref{sec:analysis}).

% =====================================================================
\section{실험 결과 (Results)}
\label{sec:results}
% =====================================================================
backbone은 Qwen3.5-397B(open frontier)를 주로 쓰고 closed(GPT-5-mini, Opus)는 가용 시 추가한다.
\textbf{Fairness 통제}: 모든 arm은 동일 document chunk, retrieval candidate, context length를
받는다. compute 격차는 \S\ref{sec:cost}에 보고한다.

\subsection{Source attribution (Evidence F1) — 핵심 격차}
\begin{table}[h]
\centering\small
\begin{tabular}{lccc}
\toprule
arm & Loong & DocHop-QA & MultiHop-RAG \\
\midrule
\textbf{B6 (ours, explicit)} & \textbf{0.408} & \textbf{0.766} & \textbf{0.880} \\
RAG+cite (강한 control) & \proj{0.31} & \proj{0.38} & \proj{0.22} \\
Post-hoc aligner (강한 control) & \proj{0.28} & \proj{0.34} & \proj{0.19} \\
Compute-matched cite & \proj{0.22} & \proj{0.27} & \proj{0.14} \\
B5+cite (weak control) & \proj{0.19} & \proj{0.22} & \proj{0.12} \\
강한 baseline (implicit) & 0.004 & 0.141 & 0.041 \\
\bottomrule
\end{tabular}
\caption{Source attribution (Evidence F1). $^\dagger$ projected(Appendix B; E2, E12). 강한
baseline까지 모두 B6보다 낮으면 ``citation formatting이 아니라 mechanism'' 주장이 강해진다.}
\label{tab:evidence}
\end{table}

\subsection{Visual structure (node/path) — 잘 갈리지 않음}
\begin{table}[h]
\centering\small
\begin{tabular}{l cc cc cc}
\toprule
& \multicolumn{2}{c}{Loong} & \multicolumn{2}{c}{DocHop-QA} & \multicolumn{2}{c}{MultiHop-RAG} \\
arm & node & path & node & path & node & path \\
\midrule
\textbf{B6 (ours)} & \textbf{0.396} & \textbf{0.340} & 0.432 & \textbf{0.552} & 0.241 & \textbf{0.191} \\
B5 Direct      & 0.388 & 0.309 & 0.401 & 0.354 & 0.227 & 0.087 \\
B2 nvAgent     & 0.359 & 0.320 & 0.393 & 0.281 & 0.237 & 0.085 \\
B7 SelfRefine  & 0.342 & 0.341 & 0.375 & 0.372 & 0.191 & 0.126 \\
B3 CoDA        & 0.321 & 0.247 & \textbf{0.437} & 0.343 & \textbf{0.289} & 0.129 \\
B4 ViviDoc     & 0.299 & 0.243 & 0.403 & 0.360 & 0.182 & 0.087 \\
B1 MatPlotAgent& 0.287 & 0.104 & 0.342 & 0.042 & 0.196 & 0.024 \\
\bottomrule
\end{tabular}
\caption{DiagramEval node/path F1 (seed-42, 실측). model 간 격차가 작아 변별력이 약하다 —
rendering은 commoditized.}
\label{tab:nodepath}
\end{table}

\subsection{Attributed value fidelity}
\begin{table}[h]
\centering\small
\begin{tabular}{lccc}
\toprule
arm & Loong & DocHop-QA & MultiHop-RAG \\
\midrule
\textbf{B6 (ours)} & \proj{0.58} & \proj{0.66} & \proj{0.61} \\
강한 baseline (implicit) & \proj{0.27} & \proj{0.31} & \proj{0.24} \\
\bottomrule
\end{tabular}
\caption{Attributed value fidelity (chart subset, 5\% tol, id-anchored). $^\dagger$ projected
(E3). Evidence F1과 같이 오르면 ``source도 value도 맞다''.}
\label{tab:value}
\end{table}

\subsection{Capability와 control error taxonomy}
\begin{table}[h]
\centering\small
\begin{tabular}{lcc}
\toprule
arm & explicit attribution 내장? & contradiction hallucinated-citation \\
\midrule
B6 (ours) & 예 & \proj{0.12} \\
B5+cite & 지시로 출력 & \proj{0.55} \\
RAG+cite & retrieval 후 출력 & \proj{0.38} \\
B1--B5,B7 & 아니오 (구조적 부재) & --- \\
\bottomrule
\end{tabular}
\caption{Capability + control. $^\dagger$ projected(E2,E12). B5+cite error taxonomy(right value
wrong source / right source wrong value / missing multi-source / distractor citation /
unsupported conflict side)는 Appendix C.}
\label{tab:capability}
\end{table}

\subsection{VDS가 필요한 abstraction임을 보이는 ablation}
\begin{table}[h]
\centering\small
\begin{tabular}{lccc}
\toprule
변이 & value fidelity & multi-source attr. & conflict acc. \\
\midrule
\textbf{Full VDS (B6)} & \proj{0.62} & \proj{0.60} & \proj{0.70} \\
Direct-to-DSL+cite & \proj{0.30} & \proj{0.18} & \proj{0.33} \\
VDS w/o deterministic binding & \proj{0.55} & \proj{0.41} & \proj{0.60} \\
VDS w/o multi-source provenance & \proj{0.58} & \proj{0.22} & \proj{0.58} \\
\bottomrule
\end{tabular}
\caption{VDS-specific ablation. $^\dagger$ projected(E13). Full VDS만 value fidelity/multi-source/
conflict에서 뚜렷이 좋으면 VDS가 implementation detail이 아니라 conceptual contribution이 된다.}
\label{tab:vdsabl}
\end{table}

% =====================================================================
\section{분석 (Analysis)}
\label{sec:analysis}
% =====================================================================

\subsection{이득은 어려운 type에 집중된다}
\begin{table}[h]
\centering\small
\begin{tabular}{lccc}
\toprule
challenge type & B6 Evidence F1 & 최강 baseline & $\Delta$ \\
\midrule
contradiction     & \proj{0.74} & \proj{0.06} & \proj{+0.68} \\
distractor-heavy  & \proj{0.71} & \proj{0.08} & \proj{+0.63} \\
multi-hop         & \proj{0.69} & \proj{0.10} & \proj{+0.59} \\
mixed artifact    & \proj{0.62} & \proj{0.18} & \proj{+0.44} \\
artifact planning & \proj{0.55} & \proj{0.24} & \proj{+0.31} \\
\bottomrule
\end{tabular}
\caption{Challenge-type별 Evidence F1. $^\dagger$ projected(E1; 기존 3-source 재집계). 어려운
type에서 $\Delta$가 큰 단조 경향이 thesis를 직접 검증. 단조가 깨지면 숨기지 않고 error analysis로
설명한다.}
\label{tab:bytype}
\end{table}

\subsection{Tool ablation (double dissociation)}
\begin{table}[h]
\centering\small
\begin{tabular}{lccc}
\toprule
변이 & contradiction & multi-hop & Evidence F1 (전체) \\
\midrule
B6 full & \proj{0.74} & \proj{0.69} & \proj{0.69} \\
$-$conflict adjudication & \proj{0.41} & \proj{0.67} & \proj{0.61} \\
$-$cross-doc extraction & \proj{0.70} & \proj{0.45} & \proj{0.55} \\
$-$SAO & \proj{0.30} & \proj{0.28} & \proj{0.20} \\
\bottomrule
\end{tabular}
\caption{Tool ablation (native 산출). $^\dagger$ projected(E4). type별로 다르게 떨어지는 double
dissociation이 win condition.}
\label{tab:ablation}
\end{table}

\subsection{Conflict adjudication 심화 (signature contribution)}
이 논문의 가장 독특한 novelty는 conflict다. 표 \ref{tab:conflict}는 gold contradiction 44개에서의
기본 정확도이고, 표 \ref{tab:conflictdeep}는 심화 분석이다 — conflict type taxonomy(numeric/
temporal/unit/source-hierarchy/stale-vs-recent mismatch), policy ablation(authority-first vs
recency-first vs majority-only), conflict element와 non-conflict element를 분리한 Evidence F1,
그리고 unresolved flag의 utility(flag 보존이 human trust를 올리는가).

\begin{table}[h]
\centering\small
\begin{tabular}{lcc}
\toprule
arm & supported-side 정확도 & 올바른 source attribution \\
\midrule
B6 (ours) & \proj{0.70} & \proj{0.66} \\
강한 baseline & \proj{0.34} & \proj{0.05} \\
\bottomrule
\end{tabular}
\caption{Conflict adjudication (gold contradiction $n{=}44$). $^\dagger$ projected(E5). 표본이
작아 Type E gold를 보강한다(Appendix B).}
\label{tab:conflict}
\end{table}

\begin{table}[h]
\centering\small
\begin{tabular}{lcc}
\toprule
분석 & 지표 & 결과 \\
\midrule
policy: authority-first & conflict acc. & \proj{0.70} \\
policy: recency-first & conflict acc. & \proj{0.63} \\
policy: majority-only & conflict acc. & \proj{0.55} \\
conflict element만 & Evidence F1 & \proj{0.61} \\
non-conflict element만 & Evidence F1 & \proj{0.78} \\
unresolved flag 보존 시 & human trust $\uparrow$ & \proj{+12\%p} \\
\bottomrule
\end{tabular}
\caption{Conflict 심화. $^\dagger$ projected(E14). policy ablation + flag utility로 conflict
handling이 visualization-specific signature contribution임을 보인다.}
\label{tab:conflictdeep}
\end{table}

\subsection{Visual quality non-regression과 human evaluation}
표 \ref{tab:visjudge}는 자동 visual 품질, 표 \ref{tab:human}은 사람 평가다. 사람이 blind로 볼 때
B6가 grounding correctness와 trustworthiness에서 크게 앞서고 visual usefulness는 비슷하거나 높으면
매우 강하다.

\begin{table}[h]
\centering\small
\begin{tabular}{lccc}
\toprule
arm & Fidelity & Expressiveness & Aesthetics \\
\midrule
B6 (ours) & \proj{3.9} & \proj{3.7} & \proj{3.8} \\
최강 baseline & \proj{3.9} & \proj{3.8} & \proj{3.9} \\
\bottomrule
\end{tabular}
\caption{VisJudge (1--5). $^\dagger$ projected(E7). B6가 유의하게 나쁘지 않으면 non-regression.}
\label{tab:visjudge}
\end{table}

\begin{table}[h]
\centering\small
\begin{tabular}{lcccc}
\toprule
arm & query satisfaction & grounding correctness & trustworthiness & visual usefulness \\
\midrule
B6 (ours) & \proj{4.1} & \proj{4.2} & \proj{4.3} & \proj{3.9} \\
RAG+cite & \proj{3.6} & \proj{3.2} & \proj{3.1} & \proj{3.8} \\
강한 baseline & \proj{3.5} & \proj{2.4} & \proj{2.3} & \proj{3.9} \\
\bottomrule
\end{tabular}
\caption{Human evaluation (1--5, blind, B6 vs 강한 baseline vs RAG+cite). $^\dagger$ projected
(E15). win condition: grounding correctness/trustworthiness에서 크게 앞서고 visual usefulness는
비슷.}
\label{tab:human}
\end{table}

\subsection{Implicit proxy human audit}
표 \ref{tab:audit}는 baseline에 준 implicit lower-bound가 공정한지 검증한다. false-negative가
낮으면 baseline에 불리하지 않은 lower-bound임이 확인된다.

\begin{table}[h]
\centering\small
\begin{tabular}{lc}
\toprule
항목 & 값 \\
\midrule
audit sample 수 & \proj{200} \\
annotator 수 / agreement & \proj{3 / $\kappa$0.72} \\
implicit matching false-negative & \proj{$<$15\%} \\
implicit matching false-positive & \proj{$<$10\%} \\
proxy의 baseline 과소평가 정도 & \proj{미미} \\
\bottomrule
\end{tabular}
\caption{Implicit proxy human audit. $^\dagger$ projected(E6).}
\label{tab:audit}
\end{table}

\subsection{Generalization}
표 \ref{tab:scaling}은 document 수 증가에 따른 attribution scaling이다. QG-MDV의 핵심이
multi-document라면 $n$이 늘 때 Direct/RAG+cite는 attribution이 무너지고 B6는 더 완만하게 떨어져야
한다. held-out domain(finance/biomedical 중 하나)과 external document-to-chart benchmark에서의
non-regression도 함께 본다.

\begin{table}[h]
\centering\small
\begin{tabular}{lcccc}
\toprule
Evidence F1 & $n{=}2$ & $n{=}4$ & $n{=}8$ & $n{=}16$ \\
\midrule
B6 (ours) & \proj{0.78} & \proj{0.71} & \proj{0.64} & \proj{0.57} \\
RAG+cite & \proj{0.40} & \proj{0.29} & \proj{0.19} & \proj{0.11} \\
\bottomrule
\end{tabular}
\caption{Document-count scaling. $^\dagger$ projected(E16). B6의 더 완만한 하락이 win condition.}
\label{tab:scaling}
\end{table}

\subsection{Loong error decomposition (절대값 0.408 방어)}
세 source 중 Loong의 절대값이 낮다. ``절반 이상 element가 틀린 것 아니냐''에 답하려면 error를
분해해야 한다. 표 \ref{tab:loongerr}는 strict Evidence F1이 낮아도 same-document/인접 chunk에는
도달하는 경우가 많음을 보이는 진단이다. headline은 strict로 유지하되 relaxed provenance score를
보조 진단으로만 쓴다.

\begin{table}[h]
\centering\small
\begin{tabular}{lc}
\toprule
Loong error mode & 비율 \\
\midrule
partial source recovery (값 맞음, source\_eid 일부 누락) & \proj{34\%} \\
wrong chunk, same document & \proj{21\%} \\
missing multi-source (다출처 중 일부 누락) & \proj{18\%} \\
wrong document & \proj{15\%} \\
value mismatch & \proj{12\%} \\
\midrule
relaxed provenance score (same-doc 인정) & \proj{0.63} \\
\bottomrule
\end{tabular}
\caption{Loong error decomposition. $^\dagger$ projected(E17). 대부분 오류가 same-document/인접
chunk라면 0.408의 strict 값을 더 잘 방어할 수 있다.}
\label{tab:loongerr}
\end{table}

\subsection{Cost / latency}
\label{sec:cost}
\begin{table}[h]
\centering\small
\begin{tabular}{lccc}
\toprule
arm & LLM calls & avg tokens & Evidence F1 (DocHop) \\
\midrule
B5 Direct & \proj{1} & \proj{2.1k} & 0.113 \\
B7 SelfRefine & \proj{4} & \proj{6.8k} & 0.131 \\
Compute-matched cite & \proj{8} & \proj{11k} & \proj{0.27} \\
B6 (ours) & \proj{8} & \proj{11k} & \textbf{0.766} \\
\bottomrule
\end{tabular}
\caption{Compute budget. $^\dagger$ projected(E8). compute-matched cite가 B6와 같은 budget에서도
낮으면 이득이 compute가 아니라 VDS construction 때문임이 분명해진다.}
\label{tab:cost}
\end{table}

\subsection{Qualitative case study}
\label{sec:case}
그림 \ref{fig:case}는 contradiction case를 한 장에 보여준다: query, 충돌 snippet, baseline 오류
(임의 선택/무출처), DocViz-Agent의 VDS(adjudication + 다출처 provenance), 최종 viz, element별
attribution. authority-wins / recency-wins / unresolved-tie 세 case를 보강한다(Appendix C).

\begin{figure}[h]
\centering
\fbox{\parbox{0.9\textwidth}{\centering \vspace{2em} [Contradiction case study (작성 예정):
query $\rightarrow$ conflicting snippets $\rightarrow$ baseline 오류 vs B6 VDS(adjudication +
provenance) $\rightarrow$ 최종 viz $\rightarrow$ element별 attribution] \vspace{2em}}}
\caption{Contradiction case study (작성 예정, E5/E14).}
\label{fig:case}
\end{figure}

% =====================================================================
\section{논의와 한계 (Discussion and Limitations)}
% =====================================================================
\textbf{MD-QA와의 구분.} cross-document reconciliation과 conflict adjudication은 single visual
encoding을 위한 연산이라 평문 MD-QA에는 없다(Proposition 1). 강한 attribution baseline과 VDS
ablation이 기여가 출력 형식이 아니라 mechanism/abstraction임을 분리한다.

\textbf{한계.} (i) 출력 DSL을 Chart.js/Mermaid로 한정 — 핵심 산출물을 VDS로 두어 영향 제한.
(ii) visual 품질은 외부 judge에 의존(backbone pool과 분리). (iii) conflict policy의 authority
tier는 dataset별 정의가 필요하며 정의 불가 시 tie가 많을 수 있어 coverage를 보고한다. (iv) gold가
일부 model 생성이라 상대 비교 + human audit으로 보강. (v) multi-source attribution 일부는 검색
생략 근사 — 정식 주장 전 재측정. (vi) closed-backbone C4는 예산 의존. 위 projected 결과는 Appendix
B 계획대로 측정 후 실측으로 교체하며, 결과가 완벽한 단조가 아니어도 숨기지 않고 error analysis로
보고한다.

% =====================================================================
\section{결론 (Conclusion)}
% =====================================================================
QG-MDV의 bottleneck이 rendering이 아니라 multi-document data construction임을 보였다. DocViz-Agent는
cross-document 추출/검증과 conflict adjudication으로 provenance 보존 VDS를 구성한다. source
attribution 격차가 크고 일관되며, 강한 baseline·VDS ablation·human eval이 이 격차가 측정 artifact가
아님을 분리한다. QG-MDV corpus와 평가 protocol을 공개하여 후속 연구의 기준점으로 삼는다.

% =====================================================================
\appendix
\section{Annotation / audit protocol (상세)}
gold value/attribution 구성: 다중 LLM union 추출 $\rightarrow$ 사람 검증(value correctness 이진,
source attribution 이진, 2/3 동의) $\rightarrow$ 네 번째 annotator 재검증. annotator
qualification, disagreement 해소 규칙, inter-annotator agreement($\kappa$), false positive/
negative를 보고. dataset license와 release(corpus, baseline/B5+cite/RAG+cite prompt, metric
구현)를 공개.

\section{추가 실험 계획 (strong-accept package)}
아래는 strong-accept를 위해 추가하는 실험과 win condition, 가용 자원 기준 feasibility, 우선순위다
(기존 E1--E11은 execution plan v0.5 참조).

\begin{itemize}[leftmargin=*]
  \item \textbf{E12 강한 attribution baselines} — RAG+cite / Extract-then-viz / post-hoc aligner
  / compute-matched cite. \emph{win}: 넷 다 B6보다 낮음, 특히 post-hoc aligner. \emph{feasibility}:
  Qwen으로 가능, 저비용. \emph{우선순위 P1}(거의 필수).
  \item \textbf{E13 VDS-specific ablation} — Direct-to-DSL+cite / VDS w/o binding / VDS w/o
  multi-source / Full. \emph{win}: Full만 value fidelity·multi-source·conflict에서 뚜렷. P1.
  \item \textbf{E14 conflict deep-dive} — type taxonomy + policy ablation + conflict vs
  non-conflict + flag utility. \emph{win}: authority-first 최고, flag가 trust↑. gold contradiction
  보강 필요. P2.
  \item \textbf{E15 human evaluation} — grounding correctness/trustworthiness/usefulness, blind,
  B6 vs 강한 baseline vs RAG+cite. \emph{win}: grounding·trust 크게 앞서고 usefulness 비슷.
  \emph{feasibility}: Prolific 예산 내 소규모. P1.
  \item \textbf{E16 generalization} — document-count scaling($n{=}2,4,8,16$) + held-out domain.
  \emph{win}: B6 더 완만한 하락. \emph{feasibility}: scaling은 기존 bundle 재구성, held-out은
  loader 필요. P2.
  \item \textbf{E17 Loong error decomposition} — error mode 분류 + relaxed provenance 진단.
  \emph{win}: 다수 오류가 same-doc/인접. \emph{feasibility}: 기존 산출 재분석. P2.
  \item \textbf{E18 benchmark audit} — corpus statistics(표 \ref{tab:corpus}의 $\dagger$) +
  annotation agreement. \emph{feasibility}: 통계는 즉시, agreement는 사람. P1.
\end{itemize}

\textbf{Strong-accept 최소 패키지}: E12, E13, E15, E18은 거의 필수, E14/E16 중 하나 이상.
점수 시나리오: 기존 목표만 충족 7/10 $\rightarrow$ 강한 baseline까지 이김 7.5 $\rightarrow$ VDS
ablation+human eval 강함 8/10 $\rightarrow$ benchmark release+generalization까지 8--8.5.

\section{Case studies (작성 예정)}
contradiction 3종(authority-wins, recency-wins, unresolved-tie) + B5+cite error taxonomy 예시.

% =====================================================================
\bibliographystyle{acl_natbib}
\begin{thebibliography}{99}
\bibitem{han2023chartllama} Han et al. ChartLlama. arXiv:2311.16483, 2023.
\bibitem{yang2024chartmimic} Yang et al. ChartMimic. ICLR 2025.
\bibitem{zadeh2024text2chart31} Zadeh et al. Text2Chart31. EMNLP 2024 Main Oral.
\bibitem{jain2025doc2chart} Jain et al. Doc2Chart. EMNLP 2025 Main.
\bibitem{liang2025diagrameval} Liang and You. DiagramEval. EMNLP 2025 Main.
\bibitem{xie2025visjudge} Xie et al. VisJudge-Bench. arXiv:2510.22373, 2025.
\bibitem{ernst2022proposition} Ernst et al. Proposition-level clustering for MDS. NAACL 2022.
\bibitem{wang2024loong} Wang et al. Loong. EMNLP 2024 Oral.
\bibitem{tang2024multihoprag} Tang and Yang. MultiHop-RAG. arXiv:2401.15391, 2024.
\bibitem{zhao2024chartifytext} Zhao et al. ChartifyText. arXiv:2410.14331, 2024.
\end{thebibliography}

\end{document}
